\scanpage{10}
\deleted{\unsure{pour|par} $=$ sous-groupe cyclique d'ordre $4$, noté $\unsure{\mathfrak{D}|\mathscr{D}}^+$ -- et $S$ et $A$ sont également des torseurs sous \unsure{$\Gamma$} $\mathfrak{D}^+$.}

\deleted{groupe} d'ordre $2$ \deleted{central} qui est le centre, et dont l'élément $\ne1$ est \uline{l'antipodisme} $\underline a$.

Il y a dans $\mathfrak{D}$ exactement deux éléments \inserted{$u,u'$} d'ordre $4$, inverses l'un de l'autre, engendrant l'unique sous-groupe cyclique d'ordre $4$, noté $\mathfrak{D}^+$ -- et $S$ et $A$ sont des torseurs sous ce même groupe. \margin{c'est lui $\Gamma$} L'antipodisme est l'unique élément d'ordre $2$ de $\mathfrak{D}^+$. \deleted{Le quotient $\mathfrak{D}/\mathfrak{D}^+$ opère sur}

\deleted{$\mathfrak{D}_{ab}^+\simeq\mathfrak{D}/\{1,\underline a\}$}

Le quotient $\mathfrak{D}/V_S\simeq\deleted{\unsure{\mathrm{op}\,\{1\}|\{\pm1\}}}\mathbb{Z}/2\mathbb{Z}$ opère sur $V_S$ par l'identité sur $\{1,\underline a\}\simeq\mathbb{F}_2\subset V_S\simeq\mathbb{F}_2^C$, et par la transposition sur $V_S\setminus\mathbb{F}_2\simeq C$, et symétriquement pour $\mathfrak{D}/V_A\simeq\mathbb{Z}/2\mathbb{Z}$ opérant sur $V_A\simeq\deleted{\unsure{V}}\inserted{\mathbb{F}_2^D}$, \deleted{\unsure{essentiellement}} via \deleted{la transpo} l'opération par transposition sur $\underline D$.

Le quotient $\mathfrak{D}/\mathfrak{D}^+\simeq\deleted{\unsure{\mathrm{op}}}\mathbb{Z}/2\mathbb{Z}$ opère sur $\mathfrak{D}^+$ par l'autom \deleted{\unsure{de ren de sy|de ren de sym}} $x\longmapsto x^{-1}$, \deleted{\unsure{sur}} \inserted{en particulier} $u\longmapsto u'=u^{-1}$ i.e. l'opération induit la transposition sur l'ens $\underline\omega_Q$ des deux générateurs de \deleted{\unsure{$\mathbb{Z}$|$\mathfrak{D}$}} $\mathfrak{D}^+$. \margin{$\mathfrak{D}^+\simeq\mathbb{Z}/4\mathbb{Z}\mathbin{\wedge}_{\{\pm1\}}\underline\omega_Q$}

On a
\[
\mathfrak{D}_{ab}^+\simeq\mathfrak{D}/\{1,\underline a\},
\]
espace vect. de rg $2$ sur $\mathbb{F}_2$. Les groupes $V_S,\mathfrak{D}^+,V_A$ sont les images inverses des trois sous-groupes d'ordre $2$ de $\mathfrak{D}_{ab}^+$, donc les trois seuls sous-groupes d'indice $2$ de $\mathfrak{D}_{ab}^+$. On a une base canonique de $\mathfrak{D}_{ab}^+$
% slanted
\margin{$\mathfrak{D}$ opère sur $C$ et sur $D$ par autom intérieurs, donc sur $C\amalg D$ en fixant $C,D$, d'où $\mathfrak{D}\to\unsure{\mathfrak{S}|G}_C\times\mathfrak{S}_D\simeq\mathbb{F}_2\times\mathbb{F}_2$, et cet hom est surjectif, et a comme noyau $\{1,\underline a\}$, et définit l'isom. can. $\mathfrak{D}/\{1,\underline a\}\simeq\mathbb{F}_2\times\mathbb{F}_2$ \unsure{indiqué ci-contre|indiqué} \unsure{les éléments de|les involutions par lesquelles} $D$ \unsure{opèrent|opérant|opère} sur $C$ par \unsure{transposition|transpositions}, \unsure{ceux de|de même} $D$ opèrent sur $C$ \unsure{par des transpositions|par transposition} (i.e. $c\in C,\ d\in D\Longrightarrow[c,d]=\underline a$)}
% slanted
\margin{\deleted{\uline{NB} $\mathfrak{D}_{ab}\overset{\sim}{\longrightarrow}V_S/\unsure{\mathbb{F}_2}\times V_A/\unsure{\mathbb{F}_2}\simeq\mathbb{F}_2^C\unsure{/\mathbb{F}_2}\times\mathbb{F}_2^D\unsure{/\mathbb{F}_2}$}}
\[
\mathfrak{D}_{ab}^+\simeq\mathbb{F}_2\times\mathbb{F}_2
\]
définie par la condition que
% archive sheet number: 9
